\documentclass[a4paper,12pt]{article}

% --- Packages ---
\usepackage[utf8]{inputenc}
\usepackage[T1]{fontenc}
\usepackage[french]{babel}
\usepackage{geometry}
\geometry{margin=2.5cm}
\usepackage{graphicx}
\usepackage{xcolor}
\usepackage{tikz}
\usetikzlibrary{shapes.geometric, arrows.meta, positioning, calc, decorations.pathreplacing}
\usepackage{amsmath}
\usepackage{booktabs}
\usepackage{longtable}
\usepackage{enumitem}
\usepackage{hyperref}
\usepackage{fancyhdr}
\usepackage{titlesec}
\usepackage{tcolorbox}
\tcbuselibrary{skins, breakable}
\usepackage{colortbl}
\usepackage{array}
\usepackage{amssymb}

% --- Colors ---
\definecolor{bmsprimary}{HTML}{1B4F72}
\definecolor{bmssecondary}{HTML}{2E86C1}
\definecolor{bmsaccent}{HTML}{E74C3C}
\definecolor{bmsgreen}{HTML}{27AE60}
\definecolor{bmsorange}{HTML}{E67E22}
\definecolor{bmsgray}{HTML}{7F8C8D}
\definecolor{bmslight}{HTML}{EBF5FB}
\definecolor{phase1}{HTML}{2980B9}
\definecolor{phase2}{HTML}{27AE60}
\definecolor{phase3}{HTML}{8E44AD}
\definecolor{phase4}{HTML}{E67E22}
\definecolor{phase5}{HTML}{C0392B}
\definecolor{phase6}{HTML}{16A085}
\definecolor{phase7}{HTML}{2C3E50}
% Status colors
\definecolor{statusdone}{HTML}{1E8449}
\definecolor{statusinprog}{HTML}{D4AC0D}
\definecolor{statustodo}{HTML}{95A5A6}

% --- Header/Footer ---
\pagestyle{fancy}
\fancyhf{}
\fancyhead[L]{\small\textcolor{bmsprimary}{\textbf{BMS Intelligent -- Roadmap}}}
\fancyhead[R]{\small\textcolor{bmsgray}{Samsung SDI 60Ah Prismatic}}
\fancyfoot[C]{\thepage}
\renewcommand{\headrulewidth}{0.4pt}

% --- Title formatting ---
\titleformat{\section}{\Large\bfseries\color{bmsprimary}}{}{0em}{\thesection\quad}
\titleformat{\subsection}{\large\bfseries\color{bmssecondary}}{}{0em}{\thesubsection\quad}
\titleformat{\subsubsection}{\normalsize\bfseries\color{bmsorange}}{}{0em}{\thesubsubsection\quad}

% --- Status badges ---
\newcommand{\done}{\textcolor{statusdone}{\textbf{[\checkmark~TERMINÉ]}}}
\newcommand{\inprog}{\textcolor{statusinprog}{\textbf{[\rightarrow~EN COURS]}}}
\newcommand{\todo}{\textcolor{statustodo}{\textbf{[---~À FAIRE]}}}

% --- Custom boxes ---
\newtcolorbox{phasebox}[2][]{
  colback=#2!5,
  colframe=#2,
  fonttitle=\bfseries\large,
  title=#1,
  breakable,
  left=5pt, right=5pt
}

\newtcolorbox{deliverablebox}[1][]{
  colback=bmsgreen!5,
  colframe=bmsgreen,
  fonttitle=\bfseries,
  title=#1,
  breakable
}

\newtcolorbox{infobox}[1][]{
  colback=bmslight,
  colframe=bmsprimary,
  fonttitle=\bfseries,
  title=#1,
  breakable
}

% ===========================================================================
\begin{document}

% --- Title page ---
\begin{titlepage}
\centering
\vspace*{2cm}
{\Huge\bfseries\textcolor{bmsprimary}{Roadmap Complet et Détaillé}}\\[0.5cm]
{\LARGE\textcolor{bmssecondary}{Modélisation et Validation d'un BMS Intelligent}}\\[0.3cm]
{\Large\textcolor{bmssecondary}{pour Véhicules Électriques -- MIL/SIL/PIL}}\\[1.5cm]
{\large Plan d'Exécution Étape par Étape}\\[2cm]

\begin{tcolorbox}[colback=bmslight, colframe=bmsprimary, width=12cm]
\centering
\textbf{7 Phases}~: de l'analyse préliminaire à la livraison finale\\[5pt]
\textbf{Chimie}~: Samsung SDI 60Ah Prismatic\\
\textbf{Dataset}~: Battery and Heating Data in Real Driving Cycles\\
\textbf{Outils}~: MATLAB/Simulink, Python, Embedded Coder, Stateflow\\
\textbf{Normes}~: ISO 26262 (ASIL)
\end{tcolorbox}

\vfill
{\large\textcolor{bmsgray}{\today}}
\end{titlepage}

\tableofcontents
\newpage

% ===========================================================================
\section*{Statut du Projet (1\textsuperscript{er} avril 2026)}
\addcontentsline{toc}{section}{Statut du Projet}
% ===========================================================================

\begin{tcolorbox}[colback=bmslight, colframe=bmsprimary]
\centering
\textbf{Avancement global~: Phase 1 terminée $\bullet$ Phase 2 terminée $\bullet$ Phase 3--7 à venir}
\end{tcolorbox}

\vspace{4pt}
\begin{center}
\renewcommand{\arraystretch}{1.4}
\begin{tabular}{c p{4cm} p{5.5cm} c}
\toprule
\textbf{Phase} & \textbf{Intitulé} & \textbf{Avancement} & \textbf{Statut} \\
\midrule
1 & Analyse Préliminaire & Data cleaning, EDA, rapports complets & \textcolor{statusdone}{\textbf{Terminé}} \\
2 & Modélisation Plant & ECM~2RC, LUTs, thermique, pack~96S validés & \textcolor{statusdone}{\textbf{Terminé}} \\
3 & Développement BMS & EKF, SoH, Balancing, Fault Detection & \textcolor{statustodo}{À faire} \\
4 & Module IA & Dataset fault injection, LSTM & \textcolor{statustodo}{À faire} \\
5 & Validation MIL/SIL/PIL & Tests, code C, PIL MCU & \textcolor{statustodo}{À faire} \\
6 & Outils IA Avancés & Complexité, dead code, doc automatique & \textcolor{statustodo}{À faire} \\
7 & Intégration Finale & Assemblage, ISO~26262, livrables & \textcolor{statustodo}{À faire} \\
\bottomrule
\end{tabular}
\end{center}

\subsection*{Résultats clés obtenus}
\begin{itemize}[nosep]
  \item \textbf{Dataset~:} 70 trips (32 TripA + 38 TripB) nettoyés -- 1\,064\,082 lignes + 917\,280 lignes simulation.
  \item \textbf{ECM~2RC~:} LUTs $R_0, R_1, C_1, R_2, C_2$ sur grille $21\times7$ $(SoC \in [0,100]\%/5,\; T \in [0,30]\,^{\circ}$C/5) extraites \textbf{exclusivement sur données de mesure}.
  \item \textbf{Thermique~:} $C_{th}=1300$~J/K, $h_{cool}=0.02$~W/K, $h_{amb}=0.15$~W/K -- RMSE~T~$=0.75\,^{\circ}$C (9 trips).
  \item \textbf{Pack 96S~:} validation sur l'ensemble des 70 trips~: RMSE~V médiane~$\approx 1102$~mV, RMSE~T médiane~$\approx 1.18\,^{\circ}$C.
  \item \textbf{Domaine de validité~:} modèle validé dans l'enveloppe $SoC \in [40\%, 90\%]$, $T_{cell} \in [10, 30]\,^{\circ}$C (RMSE~V~$\approx 200$--$700$~mV). Hors-enveloppe ($SoC < 40\%$, $T < 10\,^{\circ}$C)~: précision dégradée, cohérent avec la couverture limitée du dataset.
\end{itemize}

\newpage

% ===========================================================================
\section{Vue d'Ensemble des Phases}
% ===========================================================================

\begin{center}
\begin{tikzpicture}[
  phase/.style={rectangle, draw=#1, fill=#1!10, text width=4.5cm, minimum height=1.8cm, align=center, rounded corners=5pt, font=\small},
  arrowstyle/.style={-{Stealth[length=3mm]}, thick, #1},
  node distance=0.6cm
]

\node[phase=phase1] (p1) {\textbf{Phase 1}\\Analyse Préliminaire\\et Exploration des Données};
\node[phase=phase2, below=of p1] (p2) {\textbf{Phase 2}\\Modélisation du Plant\\(Batterie + Véhicule)};
\node[phase=phase3, below=of p2] (p3) {\textbf{Phase 3}\\Développement du\\Contrôleur BMS};
\node[phase=phase4, below=of p3] (p4) {\textbf{Phase 4}\\Module Intelligence\\Artificielle};
\node[phase=phase5, below=of p4] (p5) {\textbf{Phase 5}\\Validation MIL/SIL/PIL};
\node[phase=phase6, below=of p5] (p6) {\textbf{Phase 6}\\Outils IA Avancés\\(Dead Code, Doc, Analyse)};
\node[phase=phase7, below=of p6] (p7) {\textbf{Phase 7}\\Intégration Finale\\et Livrables};

\draw[arrowstyle=phase1] (p1) -- (p2);
\draw[arrowstyle=phase2] (p2) -- (p3);
\draw[arrowstyle=phase3] (p3) -- (p4);
\draw[arrowstyle=phase4] (p4) -- (p5);
\draw[arrowstyle=phase5] (p5) -- (p6);
\draw[arrowstyle=phase6] (p6) -- (p7);

\end{tikzpicture}
\end{center}

\newpage
% ===========================================================================
% PHASE 1
% ===========================================================================
\begin{phasebox}[Phase 1 -- Analyse Préliminaire et Exploration des Données \hfill \done]{phase1}

\textit{Fondation du projet~: comprendre les données, la chimie cellulaire et préparer l'environnement de travail.}
\end{phasebox}

\subsection{Étape 1.1 -- Configuration de l'Environnement}
\begin{enumerate}[nosep]
  \item Installer MATLAB/Simulink (R2023b ou plus récent) avec les toolboxes requises~:
    \begin{itemize}[nosep]
      \item Simulink, Stateflow, Simulink Coder, Embedded Coder.
      \item Simscape Electrical (optionnel, pour modèle physique).
      \item Deep Learning Toolbox, Statistics and Machine Learning Toolbox.
    \end{itemize}
  \item Installer Python 3.10+ avec les librairies~:
    \begin{itemize}[nosep]
      \item \texttt{numpy}, \texttt{pandas}, \texttt{matplotlib}, \texttt{scikit-learn}.
      \item \texttt{tensorflow} ou \texttt{pytorch} pour les réseaux de neurones.
      \item \texttt{onnx}, \texttt{onnxruntime} pour l'export vers Simulink.
    \end{itemize}
  \item Configurer le support matériel pour PIL (ex. STM32CubeIDE, TI Code Composer Studio).
  \item Mettre en place un dépôt Git pour le versionnement du projet.
\end{enumerate}

\subsection{Étape 1.2 -- Exploration et Compréhension du Dataset}

\subsubsection{1.2.1 -- Chargement des Données de Mesure}
\begin{enumerate}[nosep]
  \item Utiliser le script \texttt{readin.m} fourni pour charger les CSV dans MATLAB.
  \item Charger les 32 fichiers TripA (été) et les 38 fichiers TripB (hiver).
  \item Documenter les différences de structure~:
    \begin{itemize}[nosep]
      \item TripA~: 28 colonnes (données partielles, problèmes de mesure signalés).
      \item TripB~: 48 colonnes (données complètes et cohérentes, avec températures de cabine détaillées).
    \end{itemize}
  \item Identifier et traiter les valeurs manquantes (NaN)~: interpolation linéaire comme suggéré dans \texttt{readin.m}.
\end{enumerate}

\subsubsection{1.2.2 -- Chargement des Données de Simulation}
\begin{enumerate}[nosep]
  \item Charger les 24 fichiers CSV (ou MAT) de simulation.
  \item Organiser par condition~:
    \begin{itemize}[nosep]
      \item 3 températures ambiantes~: $-10\,^{\circ}$C (\texttt{Neg10}), $0\,^{\circ}$C (\texttt{0}), $+10\,^{\circ}$C (\texttt{Pos10}).
      \item 4 scénarios de conduite~: Standard, Charging, Heating, Mixed.
      \item 2 bus de communication~: CAN, LIN.
    \end{itemize}
  \item Identifier les 43 signaux disponibles~: Environment, Longitudinal Vehicle, Drive Train, Battery (Pack Voltage, Battery Current, SOC, SOH, Cell Temperature), Consumption, High Voltage Heater, Auxiliaries.
\end{enumerate}

\subsubsection{1.2.3 -- Analyse Exploratoire des Données (EDA)}
\begin{enumerate}[nosep]
  \item Tracer les profils temporels pour chaque trip~:
    \begin{itemize}[nosep]
      \item Vitesse vs. temps, courant batterie vs. temps, SoC vs. temps.
      \item Température batterie vs. température ambiante.
    \end{itemize}
  \item Analyse statistique~:
    \begin{itemize}[nosep]
      \item Distribution des courants (charge vs. décharge).
      \item Plage de SoC utilisée dans les trips réels.
      \item Corrélation température ambiante / performance batterie.
    \end{itemize}
  \item Comparer mesures réelles (TripB) vs. simulation (scénarios correspondants).
  \item Identifier les phases clés~: accélération, freinage régénératif, charge, repos, chauffage.
\end{enumerate}

\subsubsection{1.2.4 -- Étude de la Cellule Samsung SDI 60Ah}
\begin{enumerate}[nosep]
  \item Rechercher la fiche technique (datasheet) de la cellule Samsung SDI 60Ah prismatic.
  \item Documenter les paramètres clés~:
    \begin{itemize}[nosep]
      \item Tension nominale ($\approx 3.7$~V), plage de tension ($2.7$--$4.2$~V).
      \item Courant max continu et crête.
      \item Plage de température opérationnelle.
      \item Résistance interne typique.
    \end{itemize}
  \item Définir la configuration du pack~: nombre de cellules en série/parallèle (ex. 96S1P pour $\approx 355$~V).
\end{enumerate}

\begin{deliverablebox}[Livrables Phase 1 \hfill \done]
\begin{itemize}[nosep]
  \item \done~Environnement de développement opérationnel.
  \item \done~Rapport d'exploration des données (EDA) -- \texttt{ai/reports/eda\_report.md}.
  \item \done~Rapport qualité des données -- \texttt{ai/reports/data\_quality\_report.md}.
  \item \done~Rapport nettoyage -- \texttt{ai/reports/data\_cleaning\_report.md}.
  \item \done~Données nettoyées -- \texttt{data/cleaned/} (mesures + simulation).
\end{itemize}
\end{deliverablebox}

\newpage
% ===========================================================================
% PHASE 2
% ===========================================================================
\begin{phasebox}[Phase 2 -- Modélisation du Plant (Batterie + Véhicule) \hfill \done]{phase2}

\textit{Construire le modèle de la plante (battery pack + environnement véhicule) dans Simulink.}
\end{phasebox}

\subsection{Étape 2.1 -- Modèle Électrochimique de la Cellule (ECM 2RC)}

\subsubsection{2.1.1 -- Extraction de la Courbe OCV(SoC)}
\begin{enumerate}[nosep]
  \item Identifier les phases de repos prolongé dans les données (courant $\approx 0$~A pendant $>30$~min).
  \item Extraire les paires $(SoC, V_{OCV})$ à l'état d'équilibre.
  \item Ajuster un polynôme ou une table de correspondance (LUT) pour $OCV = f(SoC)$.
  \item Valider sur les phases de repos des données de mesure (courant $\approx 0$~A, TripA/TripB).
\end{enumerate}

\subsubsection{2.1.2 -- Identification des Paramètres du Circuit}
\begin{enumerate}[nosep]
  \item Identifier $R_0$ (résistance ohmique)~: mesurer le saut instantané de tension lors d'un échelon de courant.
  \item Identifier $R_1, C_1$ (première boucle RC)~: ajustement exponentiel sur le transitoire rapide.
  \item Identifier $R_2, C_2$ (deuxième boucle RC)~: ajustement sur le transitoire lent.
  \item Construire des LUT 2D~: $R_0(SoC, T)$, $R_1(SoC, T)$, $C_1(SoC, T)$, $R_2(SoC, T)$, $C_2(SoC, T)$.
  \item LUTs construites \textbf{exclusivement sur les données de mesure}~: plage de température couverte $0$--$30\,^{\circ}$C, 7 points ($0, 5, 10, 15, 20, 25, 30\,^{\circ}$C).
\end{enumerate}

\subsubsection{2.1.3 -- Implémentation Simulink du Modèle Cellule}
\begin{enumerate}[nosep]
  \item Créer un sous-système \texttt{Cell\_ECM\_2RC} dans Simulink.
  \item Entrées~: courant $I(t)$, température $T$.
  \item Sorties~: tension terminale $V_{terminal}(t)$, SoC interne (Coulomb counting).
  \item Intégrer les LUT de paramètres avec interpolation bilinéaire.
  \item Valider la réponse de la cellule vs. données mesurées.
\end{enumerate}

\subsection{Étape 2.2 -- Modèle du Pack Batterie}
\begin{enumerate}[nosep]
  \item Assembler $N$ cellules en série (96S) pour former le pack.
  \item Introduire des dispersions de paramètres entre cellules ($\pm 2\%$ sur capacité, $\pm 5\%$ sur $R_0$) pour simuler le besoin d'équilibrage.
  \item Calculer la tension pack~: $V_{pack} = \sum_{i=1}^{96} V_{cell,i}$.
  \item Calculer le courant pack (même courant en série).
\end{enumerate}

\subsection{Étape 2.3 -- Modèle Thermique}
\begin{enumerate}[nosep]
  \item Modèle thermique par cellule~: $C_{th} \frac{dT}{dt} = P_{loss} - \frac{T - T_{amb}}{R_{th}}$.
  \item Puissance dissipée~: $P_{loss} = I^2 \cdot R_0 + I^2 \cdot R_1 \cdot (1 - e^{-t/\tau_1}) + \ldots$
  \item Intégrer les données de chauffage du dataset (Heating Power CAN/LIN).
  \item Valider contre les profils de température mesurés dans TripB.
\end{enumerate}

\subsection{Étape 2.4 -- Modèle du Circuit de Refroidissement (\texttt{Coolant\_loop})}
\begin{enumerate}[nosep]
  \item Modéliser la dynamique thermique du liquide de refroidissement~:
    \[ C_{cool}\,\frac{dT_{cool}}{dt} = \dot{m}_{cool}\,c_p\,(T_{cool,in} - T_{cool}) + Q_{cell \to cool} \]
  \item Entrée~: température d'entrée $T_{cool,in}(t)$ issue du signal \texttt{coolant\_temp\_inlet\_c} (disponible depuis TripA23 et tous les TripB).
  \item Sortie~: $T_{cool}(t)$ injectée dans l'équation thermique de chaque cellule.
  \item Créer le sous-système Simulink \texttt{Coolant\_loop} et le coupler au modèle thermique.
  \item Valider contre les mesures de température du liquide disponibles dans TripA23+ et TripB.
\end{enumerate}

\subsection{Étape 2.5 -- Modèle Véhicule Simplifié}
\begin{enumerate}[nosep]
  \item Entrées du profil de conduite issu du dataset~: vitesse, élévation, couple moteur.
  \item Modèle de demande de puissance~: \( P_{elec} = f(v, a, \theta, m) \).
  \item Calcul du courant batterie demandé~: $I_{bat} = P_{elec} / V_{pack}$.
  \item Alternative~: injecter directement le courant batterie mesuré depuis les CSV.
\end{enumerate}

\begin{deliverablebox}[Livrables Phase 2 \hfill \done]
\begin{itemize}[nosep]
  \item \done~Modèle Simulink \texttt{Cell\_ECM\_2RC} validé -- \texttt{plant\_model/models/cell\_ecm\_2rc.slx}.
  \item \done~LUTs $R_0, R_1, C_1, R_2, C_2 = f(SoC, T)$ -- grille $21\times7$ -- \texttt{plant\_model/data/*.csv}.
  \item \done~Extraction OCV$(SoC)$ -- \texttt{plant\_model/data/ocv\_lut.csv}.
  \item \done~Modèle du pack 96S avec dispersions -- \texttt{plant\_model/models/pack\_96s.slx}.
  \item \done~Modèle thermique calibré~: $C_{th}=1300$~J/K, RMSE~$=0.75\,^{\circ}$C.
  \item \done~Modèle \texttt{Coolant\_loop} (circuit de refroidissement) couplé au modèle thermique.
  \item \done~Validation pack (données de mesure uniquement)~: RMSE~V médiane~$\approx 1102$~mV, RMSE~T médiane~$\approx 1.18\,^{\circ}$C.
  \item \done~Scripts -- \texttt{extract\_ecm\_params.py}, \texttt{calibrate\_thermal.m}, \texttt{validate\_pack\_model.m}.
\end{itemize}
\end{deliverablebox}

\newpage
% ===========================================================================
% PHASE 3
% ===========================================================================
\begin{phasebox}[Phase 3 -- Développement du Contrôleur BMS \hfill \todo]{phase3}

\textit{Implémenter les fonctions principales du BMS dans Simulink~: SoC, SoH, équilibrage, détection de défauts, puis re-valider le système complet sur les données de mesure.}

\smallskip
\textbf{Architecture~:} Le BMS adopte une structure \textbf{Master--Slave} (1~Master / 8~Slaves, un slave par module~12S). Cette architecture est prise en compte dans l'organisation des sous-systèmes Simulink.
\end{phasebox}

\subsection{Étape 3.1 -- Estimateur SoC (Extended Kalman Filter)}

\subsubsection{3.1.1 -- Implémentation de l'EKF}
\begin{enumerate}[nosep]
  \item Définir le vecteur d'état~: $\mathbf{x} = [SoC,\; V_{RC1},\; V_{RC2}]^T$.
  \item Implémenter l'équation d'état (prédiction)~:
    \[
    \begin{bmatrix} SoC_{k+1} \\ V_{RC1,k+1} \\ V_{RC2,k+1} \end{bmatrix}
    = f\left(\begin{bmatrix} SoC_k \\ V_{RC1,k} \\ V_{RC2,k} \end{bmatrix}, I_k \right)
    \]
  \item Implémenter l'équation de mesure~: $V_{terminal} = OCV(SoC) - I \cdot R_0 - V_{RC1} - V_{RC2}$.
  \item Calculer les matrices Jacobiennes $\mathbf{F}_k$ et $\mathbf{H}_k$.
  \item Régler les matrices de covariance $\mathbf{Q}$ (bruit de processus) et $R$ (bruit de mesure).
  \item Créer le sous-système Simulink \texttt{SoC\_Estimator\_EKF}.
\end{enumerate}

\subsubsection{3.1.2 -- Validation du SoC}
\begin{enumerate}[nosep]
  \item Tester sur les profils TripB (données complètes) en injectant $I(t)$ et $V(t)$.
  \item Comparer $\hat{SoC}$ estimé vs. SoC mesuré du dataset.
  \item Critère~: erreur RMS $<2\%$ sur l'ensemble des trips.
  \item Tester la robustesse~: initialisation avec erreur de SoC initiale ($\pm 10\%$).
  \item Valider sur les différentes conditions de température des données de mesure ($0$--$30\,^{\circ}$C).
\end{enumerate}

\subsection{Étape 3.2 -- Estimateur SoH}
\begin{enumerate}[nosep]
  \item Implémenter le suivi de $R_0$ par estimation récursive (RLS ou EKF étendu).
  \item Calculer $SoH_R = R_{0,init} / R_{0,actuel} \times 100\%$.
  \item Implémenter l'estimation de capacité résiduelle via Coulomb counting sur cycles complets.
  \item Suivre l'évolution de $R_0$ sur l'ensemble des trips de mesure pour estimer le vieillissement relatif.
  \item Créer le sous-système Simulink \texttt{SoH\_Estimator}.
\end{enumerate}

\subsection{Étape 3.3 -- Équilibrage Cellulaire}
\begin{enumerate}[nosep]
  \item Implémenter l'algorithme d'équilibrage passif dans Simulink~:
    \begin{itemize}[nosep]
      \item Calculer la tension moyenne des cellules~: $\bar{V} = \frac{1}{N}\sum V_{cell,i}$.
      \item Identifier les cellules dont $V_{cell,i} - \bar{V} > \Delta V_{seuil}$ (20~mV).
      \item Activer la résistance de décharge (shunt) pour ces cellules.
    \end{itemize}
  \item Implémenter la logique dans un bloc Stateflow (états~: Idle, Measuring, Balancing, Complete).
  \item Valider en simulant un pack avec dispersions de SoC initial.
  \item Vérifier la convergence des tensions cellulaires.
\end{enumerate}

\subsection{Étape 3.4 -- Détection de Défauts (Fault Detection)}
\begin{enumerate}[nosep]
  \item Créer un Stateflow \texttt{Fault\_Manager} avec les états~:
    \begin{itemize}[nosep]
      \item \texttt{Normal}~: fonctionnement nominal.
      \item \texttt{Warning}~: seuils approchés, limitation de puissance.
      \item \texttt{Fault\_Active}~: défaut confirmé, action corrective.
      \item \texttt{Critical}~: arrêt d'urgence.
    \end{itemize}
  \item Implémenter les détections~:
    \begin{itemize}[nosep]
      \item Surtension ($V_{cell} > 4.25$~V), sous-tension ($V_{cell} < 2.7$~V).
      \item Surcourant ($|I| > 180$~A, soit 3C).
      \item Surchauffe ($T > 55\,^{\circ}$C), sous-température ($T < -20\,^{\circ}$C).
      \item Défaut capteur (signal hors plage, NaN, dérivée impossible).
    \end{itemize}
  \item Implémenter la temporisation (debounce)~: confirmation du défaut après $N$ échantillons consécutifs.
  \item Définir les actions par niveau ASIL (cf. table ISO~26262 du workflow).
  \item Associer un DTC (Diagnostic Trouble Code) à chaque défaut.
\end{enumerate}

\subsection{Étape 3.5 -- Gestion Thermique}
\begin{enumerate}[nosep]
  \item Implémenter un contrôleur PI pour la régulation de température.
  \item Stratégie de chauffage~: activer si $T_{cell} < 5\,^{\circ}$C (basé sur les données TripB hiver).
  \item Stratégie de refroidissement~: activer si $T_{cell} > 40\,^{\circ}$C.
  \item Intégrer les signaux \texttt{Heating Power CAN/LIN} et \texttt{Coolant Temperature} du dataset.
  \item Créer le sous-système \texttt{Thermal\_Controller}.
\end{enumerate}

\subsection{Étape 3.6 -- Interface CAN}
\begin{enumerate}[nosep]
  \item Définir les messages CAN du BMS~:
    \begin{itemize}[nosep]
      \item \texttt{BMS\_Status}~: SoC, SoH, état du BMS.
      \item \texttt{BMS\_Faults}~: DTC actifs, niveau de sévérité.
      \item \texttt{BMS\_Limits}~: courant max charge/décharge, puissance disponible.
      \item \texttt{BMS\_Thermal}~: températures, état chauffage/refroidissement.
    \end{itemize}
  \item Modéliser l'interface CAN dans Simulink (Vehicle Network Toolbox ou bloc simplifié).
\end{enumerate}

\begin{deliverablebox}[Livrables Phase 3]
\begin{itemize}[nosep]
  \item Sous-système \texttt{SoC\_Estimator\_EKF} validé (erreur $<2\%$).
  \item Sous-système \texttt{SoH\_Estimator} validé.
  \item Stateflow \texttt{Cell\_Balancing} fonctionnel.
  \item Stateflow \texttt{Fault\_Manager} avec tous les scénarios de défaut.
  \item \texttt{Thermal\_Controller} validé sur données hiver.
  \item Interface CAN définie.
  \item Rapport de re-validation du BMS complet (Plant + Contrôleur) sur données de mesure.
\end{itemize}
\end{deliverablebox}

\newpage
% ===========================================================================
% PHASE 4
% ===========================================================================
\begin{phasebox}[Phase 4 -- Module Intelligence Artificielle \hfill \todo]{phase4}

\textit{Développer et entraîner le modèle IA de prédiction de défauts, et préparer le dataset d'entraînement.}
\end{phasebox}

\subsection{Étape 4.1 -- Préparation du Dataset IA}

\subsubsection{4.1.1 -- Extraction des Features depuis le Dataset Réel}
\begin{enumerate}[nosep]
  \item Extraire des fenêtres temporelles glissantes (ex. 100 pas = 10~s à 10~Hz) à partir des 70 trips.
  \item Features par fenêtre~:
    \begin{itemize}[nosep]
      \item Séries temporelles~: $V(t)$, $I(t)$, $T(t)$, $SoC(t)$.
      \item Features statistiques~: moyenne, écart-type, min, max, gradient.
      \item Features fréquentielles (optionnel)~: FFT du courant.
    \end{itemize}
  \item Label~: \texttt{Normal} pour toutes les données réelles (pas de défaut observé).
\end{enumerate}

\subsubsection{4.1.2 -- Génération du Dataset de Défauts Simulés (Fault Injection)}
\begin{enumerate}[nosep]
  \item \textbf{Surtension}~: ajouter un offset progressif ($+0.2$ à $+0.8$~V) sur $V_{cell}$.
  \item \textbf{Sous-tension}~: soustraire un offset ($-0.3$ à $-1.0$~V) ou accélérer la décharge.
  \item \textbf{Surcourant}~: multiplier $I(t)$ par un facteur ($1.5\times$ à $4\times$).
  \item \textbf{Surchauffe}~: injecter une rampe de température anormale ($+2\,^{\circ}$C/min).
  \item \textbf{Défaut capteur}~: remplacer des segments par une valeur constante, injecter du bruit fort, ou introduire des sauts.
  \item \textbf{Court-circuit interne}~: chute brutale de tension ($-1$~V en 5~s) + élévation rapide de température.
  \item Appliquer ces injections sur les données réelles et de simulation pour générer des versions fautées.
  \item Label~: \texttt{Overvoltage}, \texttt{Undervoltage}, \texttt{Overcurrent}, \texttt{Overtemperature}, \texttt{SensorFault}, \texttt{InternalShortCircuit}.
\end{enumerate}

\subsubsection{4.1.3 -- Constitution du Dataset Final}
\begin{enumerate}[nosep]
  \item Assembler données normales + données fautées.
  \item Équilibrer les classes (SMOTE, undersampling, ou augmentation).
  \item Normaliser (Min-Max ou Z-score) par feature.
  \item Split~: 70\% entraînement, 15\% validation, 15\% test.
  \item Sauvegarder en format compatible (HDF5, NumPy, ou MAT).
\end{enumerate}

\subsection{Étape 4.2 -- Conception et Entraînement du Modèle}

\subsubsection{4.2.1 -- Architecture du Réseau}
\begin{enumerate}[nosep]
  \item Architecture proposée (LSTM)~:
    \begin{itemize}[nosep]
      \item Input~: $(batch, timesteps, features)$ = $(batch, 100, 4)$.
      \item LSTM~: 128 unités, return sequences = False.
      \item Dense~: 64 neurones, ReLU.
      \item Dense~: 32 neurones, ReLU.
      \item Dropout~: 0.3 (régularisation).
      \item Output~: 6 neurones (classes), Softmax.
    \end{itemize}
  \item Alternative~: 1D-CNN + LSTM hybride pour capturer patterns courts et longs.
\end{enumerate}

\subsubsection{4.2.2 -- Entraînement}
\begin{enumerate}[nosep]
  \item Optimizer~: Adam ($lr = 10^{-3}$, decay cosine).
  \item Loss~: Categorical cross-entropy (ou focal loss si classes déséquilibrées).
  \item Callbacks~: Early stopping (patience=10), Model checkpoint, ReduceLROnPlateau.
  \item Epochs max~: 200, batch size~: 64.
  \item Métriques~: accuracy, precision, recall, F1-score par classe, matrice de confusion.
\end{enumerate}

\subsubsection{4.2.3 -- Évaluation et Validation}
\begin{enumerate}[nosep]
  \item Évaluer sur le set de test~: F1-score $>0.95$ visé.
  \item Analyse des erreurs~: quels types de défauts sont confondus.
  \item Test de robustesse~: évaluer sur un trip jamais vu (leave-one-trip-out).
  \item Analyser la latence d'inférence (temps de prédiction par fenêtre).
\end{enumerate}

\subsection{Étape 4.3 -- Intégration dans Simulink}
\begin{enumerate}[nosep]
  \item Exporter le modèle entraîné en format ONNX.
  \item Importer dans Simulink via le bloc \texttt{Predict} (Deep Learning Toolbox) ou MATLAB Function.
  \item Alternative~: convertir en code C/C++ via TensorFlow Lite et intégrer comme S-Function.
  \item Créer le sous-système \texttt{AI\_Fault\_Predictor} dans le modèle BMS.
  \item Connecter les sorties de l'IA au \texttt{Fault\_Manager} (alerte prédictive avant le seuil).
  \item Valider l'intégration~: les prédictions IA doivent précéder la détection par seuil.
\end{enumerate}

\begin{deliverablebox}[Livrables Phase 4]
\begin{itemize}[nosep]
  \item Dataset IA complet (normal + défauts simulés), documenté.
  \item Modèle LSTM entraîné et évalué (rapport de performance).
  \item Modèle exporté en ONNX (ou TFLite).
  \item Sous-système Simulink \texttt{AI\_Fault\_Predictor} intégré et validé.
  \item Scripts Python/MATLAB d'entraînement et d'évaluation.
\end{itemize}
\end{deliverablebox}

\newpage
% ===========================================================================
% PHASE 5
% ===========================================================================
\begin{phasebox}[Phase 5 -- Validation MIL / SIL / PIL \hfill \todo]{phase5}

\textit{Appliquer le workflow de validation complet conforme aux standards automobiles.}
\end{phasebox}

\subsection{Étape 5.1 -- Validation MIL (Model-in-the-Loop)}

\subsubsection{5.1.1 -- Préparation du Test Bench MIL}
\begin{enumerate}[nosep]
  \item Créer un modèle de test (\texttt{MIL\_Test\_Bench.slx})~:
    \begin{itemize}[nosep]
      \item Plant Model (batterie + véhicule) en boucle avec le BMS Controller.
      \item Entrées~: profils de conduite issus du dataset (injection par \texttt{From Workspace} ou \texttt{Signal Builder}).
    \end{itemize}
  \item Définir la matrice de test~:
    \begin{itemize}[nosep]
      \item 70 trips réels (TripA + TripB).
      \item 12 scénarios de simulation (3 températures $\times$ 4 modes).
      \item Scénarios de défaut injectés.
    \end{itemize}
\end{enumerate}

\subsubsection{5.1.2 -- Exécution des Tests MIL}
\begin{enumerate}[nosep]
  \item Exécuter chaque scénario et enregistrer les résultats~:
    \begin{itemize}[nosep]
      \item SoC estimé vs. référence.
      \item SoH estimé vs. référence (données simulation).
      \item Détection correcte des défauts (taux de détection, faux positifs).
      \item Comportement de l'équilibrage.
      \item Prédictions IA vs. détection par seuil.
    \end{itemize}
  \item Couverture des tests~:
    \begin{itemize}[nosep]
      \item Couverture structurelle du Stateflow \texttt{Fault\_Manager} (100\% des transitions).
      \item Couverture des MC/DC (Modified Condition/Decision Coverage) -- ISO~26262 ASIL~D.
    \end{itemize}
  \item Générer un rapport MIL avec résultats pass/fail par critère.
\end{enumerate}

\subsection{Étape 5.2 -- Validation SIL (Software-in-the-Loop)}

\subsubsection{5.2.1 -- Génération du Code C}
\begin{enumerate}[nosep]
  \item Configurer le modèle BMS Controller pour la génération de code~:
    \begin{itemize}[nosep]
      \item System Target File~: \texttt{ert.tlc} (Embedded Coder).
      \item Hardware Implementation~: processeur cible (ex. ARM Cortex-M4, 32-bit, little-endian).
      \item Solver~: Fixed-step, discrete, $T_s = 10$~ms.
    \end{itemize}
  \item Générer le code C via \texttt{Ctrl+B} ou \texttt{slbuild}.
  \item Inspecter le code généré~: structure, lisibilité, conformité MISRA~C.
\end{enumerate}

\subsubsection{5.2.2 -- Exécution SIL}
\begin{enumerate}[nosep]
  \item Configurer le mode SIL dans Simulink (Block Mode ou Top-Model SIL).
  \item Exécuter les mêmes scénarios que MIL.
  \item Comparer les résultats MIL vs. SIL~:
    \begin{itemize}[nosep]
      \item Critère~: écart numérique $<10^{-6}$ (virgule flottante).
      \item Identifier les éventuelles divergences (arrondis, séquencement).
    \end{itemize}
  \item Analyser le code C généré~:
    \begin{itemize}[nosep]
      \item Utilisation mémoire estimée (RAM, ROM).
      \item Rapport de conformité MISRA~C (Model Advisor).
    \end{itemize}
\end{enumerate}

\subsection{Étape 5.3 -- Validation PIL (Processor-in-the-Loop)}

\subsubsection{5.3.1 -- Configuration du Matériel Cible}
\begin{enumerate}[nosep]
  \item Sélectionner le microcontrôleur cible (ex. STM32F4, TI TMS320, NXP S32K144).
  \item Installer le support hardware package correspondant dans MATLAB.
  \item Configurer la chaîne de compilation croisée (cross-compiler).
  \item Configurer la communication série/JTAG entre MATLAB et le MCU.
\end{enumerate}

\subsubsection{5.3.2 -- Exécution PIL}
\begin{enumerate}[nosep]
  \item Configurer le mode PIL dans Simulink.
  \item Flasher le code compilé sur le MCU.
  \item Exécuter les scénarios de test~: les stimuli sont envoyés depuis Simulink au MCU, les résultats sont renvoyés.
  \item Mesurer les performances~:
    \begin{itemize}[nosep]
      \item \textbf{Temps d'exécution}~: WCET par pas de calcul. Vérifier $T_{exec} < T_s$.
      \item \textbf{Utilisation mémoire}~: stack, heap, flash, RAM. Vérifier $<80\%$ de la capacité.
      \item \textbf{Précision numérique}~: écart SIL vs. PIL (effets de l'arithmétique du MCU).
    \end{itemize}
  \item Comparaison SIL vs. PIL~: critère de tolérance adapté à l'architecture (virgule flottante MCU ou virgule fixe).
\end{enumerate}

\subsubsection{5.3.3 -- Rapport PIL}
\begin{enumerate}[nosep]
  \item Générer le rapport PIL incluant~:
    \begin{itemize}[nosep]
      \item Résultats de comparaison numérique (graphiques superposés).
      \item Profils de timing (histogramme des temps d'exécution).
      \item Carte mémoire (memory map).
      \item Conformité aux contraintes temps réel.
    \end{itemize}
\end{enumerate}

\begin{deliverablebox}[Livrables Phase 5]
\begin{itemize}[nosep]
  \item Rapport MIL complet (résultats, couverture de test, pass/fail).
  \item Code C généré et inspecté.
  \item Rapport SIL (comparaison MIL vs. SIL, analyse mémoire, MISRA~C).
  \item Rapport PIL (comparaison SIL vs. PIL, timing, mémoire MCU).
  \item Exécutable C testé sur microcontrôleur.
\end{itemize}
\end{deliverablebox}

\newpage
% ===========================================================================
% PHASE 6
% ===========================================================================
\begin{phasebox}[Phase 6 -- Outils IA Avancés (Dead Code, Documentation, Analyse) \hfill \todo]{phase6}

\textit{Développer les outils IA auxiliaires pour l'assistance à l'ingénieur.}
\end{phasebox}

\subsection{Étape 6.1 -- Analyse de Complexité du Modèle}
\begin{enumerate}[nosep]
  \item Développer un script MATLAB (\texttt{model\_complexity\_analyzer.m})~:
    \begin{itemize}[nosep]
      \item Parcourir récursivement tous les sous-systèmes du modèle Simulink.
      \item Compter~: nombre de blocs, sous-systèmes, états Stateflow, transitions.
      \item Calculer la profondeur d'imbrication maximale.
      \item Calculer le nombre de connexions (signaux) entre sous-systèmes.
      \item Calculer la complexité cyclomatique des charts Stateflow.
    \end{itemize}
  \item Appliquer des heuristiques IA~:
    \begin{itemize}[nosep]
      \item Seuils de complexité (ex. profondeur $>5$, blocs $>100$ par sous-système).
      \item Score de complexité global (combinaison pondérée des métriques).
      \item Recommandations automatiques (refactoring, décomposition).
    \end{itemize}
  \item Générer un rapport de complexité avec visualisation (diagramme en barres, heatmap).
\end{enumerate}

\subsection{Étape 6.2 -- Détection de Dead Code}

\subsubsection{6.2.1 -- Dead Code Simulink}
\begin{enumerate}[nosep]
  \item Analyse statique~:
    \begin{itemize}[nosep]
      \item Identifier les blocs non connectés (\texttt{find\_system} avec \texttt{FindAll}).
      \item Identifier les signaux non lus (output ports non connectés).
      \item Identifier les sous-systèmes jamais activés (triggered/enabled sans source).
    \end{itemize}
  \item Utiliser le \textit{Model Advisor} de Simulink (checks existants).
  \item Compléter par une analyse de couverture de simulation (blocs jamais exécutés).
\end{enumerate}

\subsubsection{6.2.2 -- Dead Code Stateflow}
\begin{enumerate}[nosep]
  \item Analyse de couverture des états~: identifier les états jamais atteints.
  \item Analyse des transitions~: transitions avec gardes contradictoires ou jamais vraies.
  \item Analyse des actions~: actions entry/during/exit jamais exécutées.
  \item Exploiter les rapports de couverture de Simulink Coverage.
\end{enumerate}

\subsubsection{6.2.3 -- Dead Code dans le C Généré}
\begin{enumerate}[nosep]
  \item Analyser le code C généré via des outils statiques (Polyspace, cppcheck, ou script custom).
  \item Identifier les fonctions non appelées, variables non utilisées, branches mortes.
  \item Corréler avec les blocs Simulink source pour remonter l'information.
\end{enumerate}

\subsection{Étape 6.3 -- Génération Automatique de Documentation}
\begin{enumerate}[nosep]
  \item Développer un script MATLAB/Python (\texttt{doc\_generator.m} / \texttt{doc\_generator.py})~:
    \begin{itemize}[nosep]
      \item Parcourir l'arbre hiérarchique du modèle Simulink.
      \item Pour chaque sous-système, extraire~: nom, description, entrées/sorties, paramètres.
      \item Générer une fiche descriptive par sous-système.
    \end{itemize}
  \item Générer automatiquement~:
    \begin{itemize}[nosep]
      \item Table des signaux (nom, type, unité, plage, source, destination).
      \item Diagramme d'architecture (tikz ou export d'image).
      \item Table des paramètres configurables.
      \item Résumé des résultats MIL/SIL/PIL (import des rapports).
      \item Métriques de complexité et de couverture.
    \end{itemize}
  \item Format de sortie~: \LaTeX\ (compilable), PDF, ou HTML.
  \item Intégrer les résultats IA (prédiction, dead code, complexité) dans la documentation.
\end{enumerate}

\begin{deliverablebox}[Livrables Phase 6]
\begin{itemize}[nosep]
  \item Script \texttt{model\_complexity\_analyzer} + rapport de complexité.
  \item Script \texttt{dead\_code\_detector} + rapport de dead code (Simulink + Stateflow + C).
  \item Script \texttt{doc\_generator} + documentation automatique du modèle BMS.
  \item Intégration de tous les outils dans un pipeline unique.
\end{itemize}
\end{deliverablebox}

\newpage
% ===========================================================================
% PHASE 7
% ===========================================================================
\begin{phasebox}[Phase 7 -- Intégration Finale et Livrables \hfill \todo]{phase7}

\textit{Assembler tous les composants, valider l'ensemble et préparer les livrables finaux.}
\end{phasebox}

\subsection{Étape 7.1 -- Intégration du Système Complet}
\begin{enumerate}[nosep]
  \item Assembler dans un modèle Simulink unique~:
    \begin{itemize}[nosep]
      \item Plant Model (batterie + véhicule + thermique).
      \item BMS Controller (SoC, SoH, Balancing, Fault Manager, Thermal Controller, CAN).
      \item AI Module (Fault Predictor).
      \item Test Harness (MIL/SIL/PIL configs).
    \end{itemize}
  \item Vérifier la cohérence des interfaces (types de données, dimensions, pas de calcul).
  \item Exécuter un test de bout en bout avec un profil de conduite complet.
\end{enumerate}

\subsection{Étape 7.2 -- Exécution du Pipeline IA Complet}
\begin{enumerate}[nosep]
  \item Lancer l'analyseur de complexité sur le modèle final.
  \item Lancer la détection de dead code sur le modèle final.
  \item Lancer la génération automatique de documentation.
  \item Vérifier la cohérence entre les rapports IA et les résultats de validation.
\end{enumerate}

\subsection{Étape 7.3 -- Revue ISO 26262}
\begin{enumerate}[nosep]
  \item Vérifier la traçabilité exigences $\rightarrow$ implémentation $\rightarrow$ tests.
  \item Vérifier la couverture de test (MC/DC) pour les fonctions ASIL~D.
  \item Vérifier la conformité MISRA~C du code généré.
  \item Compléter le HARA et les Safety Goals.
  \item Documenter les justifications pour les anomalies ou limitations.
\end{enumerate}

\subsection{Étape 7.4 -- Préparation des Livrables Finaux}
\begin{enumerate}[nosep]
  \item Modèle Simulink complet (\texttt{BMS\_Intelligent\_Model.slx}).
  \item Code C généré et testé en PIL.
  \item Dataset IA complet (données normales + défauts simulés).
  \item Modèle IA entraîné (ONNX / TFLite).
  \item Rapports~:
    \begin{itemize}[nosep]
      \item Rapport MIL (résultats et couverture).
      \item Rapport SIL (comparaison et analyse mémoire).
      \item Rapport PIL (timing, mémoire MCU, comparaison).
      \item Rapport de complexité du modèle.
      \item Rapport de dead code.
      \item Documentation automatique complète du BMS.
    \end{itemize}
  \item Scripts et outils (analyseur, détecteur dead code, générateur de documentation).
  \item Présentation de soutenance.
\end{enumerate}

\subsection{Étape 7.5 -- Rédaction du Rapport de Stage}
\begin{enumerate}[nosep]
  \item Introduction et contexte (BMS, véhicules électriques, ISO~26262).
  \item État de l'art (BMS, estimation SoC/SoH, IA pour batteries).
  \item Méthodologie (architecture, algorithmes, workflow MIL/SIL/PIL).
  \item Résultats et discussion.
  \item Conclusion et perspectives.
\end{enumerate}

\begin{deliverablebox}[Livrables Finaux]
\begin{itemize}[nosep]
  \item Modèle BMS intelligent complet (Simulink).
  \item Exécutable C testé en PIL sur microcontrôleur.
  \item Modules IA~: prédiction de défauts, analyse modèle, dead code, documentation auto.
  \item Rapports MIL/SIL/PIL complets.
  \item Documentation automatique du système.
  \item Rapport de stage.
  \item Présentation de soutenance.
\end{itemize}
\end{deliverablebox}

\newpage
% ===========================================================================
\section{Tableau Récapitulatif des Phases}
% ===========================================================================

\begin{center}
\small
\renewcommand{\arraystretch}{1.4}
\begin{longtable}{>{\bfseries}p{0.5cm} p{3.5cm} p{5cm} p{3.5cm}}
\toprule
\textbf{\#} & \textbf{Phase} & \textbf{Activités Principales} & \textbf{Livrables Clés} \\
\midrule
\endhead
\cellcolor{phase1!15}1 & \cellcolor{phase1!15}Analyse Préliminaire & Exploration dataset, étude cellule Samsung SDI 60Ah, configuration environnement & Rapport EDA, fiche cellule \\
\midrule
\cellcolor{phase2!15}2 & \cellcolor{phase2!15}Modélisation Plant & ECM 2RC (LUTs mesure $0$--$30\,^{\circ}$C), Coolant\_loop, pack 96S, validation sur mesures & Modèle Simulink Plant validé \\
\midrule
\cellcolor{phase3!15}3 & \cellcolor{phase3!15}Développement BMS & Architecture Master-Slave (bref), SoC (EKF), SoH, Balancing, Fault Detection, Thermal, CAN, re-validation & Contrôleur BMS complet \\
\midrule
\cellcolor{phase4!15}4 & \cellcolor{phase4!15}Module IA & Dataset fault injection, LSTM entraîné, intégration Simulink & Modèle IA + sous-système Simulink \\
\midrule
\cellcolor{phase5!15}5 & \cellcolor{phase5!15}Validation MIL/SIL/PIL & Tests MIL, génération code C, SIL, déploiement PIL & Rapports MIL/SIL/PIL, exe C \\
\midrule
\cellcolor{phase6!15}6 & \cellcolor{phase6!15}Outils IA Avancés & Complexité modèle, dead code, documentation automatique & Scripts + rapports IA \\
\midrule
\cellcolor{phase7!15}7 & \cellcolor{phase7!15}Intégration Finale & Assemblage, revue ISO~26262, livrables, rapport de stage & Système complet + rapport \\
\bottomrule
\end{longtable}
\end{center}

% ===========================================================================
\section{Dépendances entre Phases}
% ===========================================================================

\begin{center}
\begin{tikzpicture}[
  node distance=1.5cm,
  pnode/.style={rectangle, draw=#1, fill=#1!15, text width=2.5cm, minimum height=1cm, align=center, rounded corners=3pt, font=\small\bfseries},
  dep/.style={-{Stealth[length=3mm]}, thick, #1}
]

\node[pnode=phase1] (p1) {Phase 1\\Analyse};
\node[pnode=phase2, right=2cm of p1] (p2) {Phase 2\\Plant};
\node[pnode=phase3, right=2cm of p2] (p3) {Phase 3\\BMS};
\node[pnode=phase4, below=1.5cm of p3] (p4) {Phase 4\\IA};
\node[pnode=phase5, left=2cm of p4] (p5) {Phase 5\\MIL/SIL/PIL};
\node[pnode=phase6, left=2cm of p5] (p6) {Phase 6\\Outils IA};
\node[pnode=phase7, below=1.5cm of p5] (p7) {Phase 7\\Intégration};

\draw[dep=phase1] (p1) -- (p2);
\draw[dep=phase2] (p2) -- (p3);
\draw[dep=phase3] (p3) -- (p4);
\draw[dep=phase3] (p3) |- (p5);
\draw[dep=phase4] (p4) -- (p5);
\draw[dep=phase5] (p5) -- (p6);
\draw[dep=phase5] (p5) -- (p7);
\draw[dep=phase6] (p6) |- (p7);

\end{tikzpicture}
\end{center}

\textbf{Note}~: Les phases 4 (IA) et la préparation de la phase 5 (configuration MIL) peuvent être partiellement parallélisées une fois la phase 3 terminée.

% ===========================================================================
\section{Outils et Technologies}
% ===========================================================================

\begin{center}
\small
\begin{tabular}{lll}
\toprule
\textbf{Catégorie} & \textbf{Outil} & \textbf{Usage} \\
\midrule
Modélisation & MATLAB/Simulink & Modèle BMS + Plant \\
Machine d'état & Stateflow & Fault Manager, Balancing \\
Génération de code & Simulink Coder / Embedded Coder & Code C pour SIL/PIL \\
IA (entraînement) & Python (TensorFlow/PyTorch) & LSTM prédiction défauts \\
IA (intégration) & ONNX / TFLite / MATLAB DL Toolbox & Import dans Simulink \\
Analyse statique & Polyspace / cppcheck / Model Advisor & MISRA~C, dead code \\
Couverture de test & Simulink Coverage & MC/DC, couverture Stateflow \\
MCU cible & STM32 / TI C2000 / NXP S32K & Platform PIL \\
Versionnement & Git & Gestion du code source \\
Documentation & \LaTeX\ / Python scripts & Rapports automatiques \\
Données & Dataset CSV/MAT & Stimuli et validation \\
\bottomrule
\end{tabular}
\end{center}

\newpage
% ===========================================================================
\section{Critères de Réussite}
% ===========================================================================

\begin{center}
\begin{tabular}{p{5cm}p{7cm}}
\toprule
\textbf{Critère} & \textbf{Objectif} \\
\midrule
Précision SoC & Erreur RMS $<2\%$ sur tous les profils de conduite \\
Précision SoH & Erreur $<5\%$ par rapport à la référence simulation \\
Détection de défauts & Taux de détection $>99\%$, faux positifs $<1\%$ \\
Prédiction IA & F1-score $>0.95$ sur le set de test \\
Comparaison MIL vs. SIL & Écart numérique $<10^{-6}$ \\
Performance PIL & $T_{exec} < T_s = 10$~ms, mémoire $<80\%$ MCU \\
Couverture de test & $>95\%$ couverture structurelle, 100\% transitions Stateflow \\
Conformité MISRA~C & 0 violation de règles obligatoires \\
Documentation & Génération automatique complète et à jour \\
\bottomrule
\end{tabular}
\end{center}

% ===========================================================================
\section{Risques et Mitigations}
% ===========================================================================

\begin{center}
\small
\begin{tabular}{p{4cm}p{2cm}p{6cm}}
\toprule
\textbf{Risque} & \textbf{Impact} & \textbf{Mitigation} \\
\midrule
Données insuffisantes pour l'IA & Élevé & Augmentation par fault injection + données simulation \\
Précision EKF insuffisante & Moyen & Ajuster $\mathbf{Q}/R$, passer à UKF si nécessaire \\
MCU mémoire insuffisante & Moyen & Optimiser le code, réduire la taille du modèle IA \\
Latence IA trop élevée & Moyen & Quantification du modèle, réduction architecture \\
Divergence MIL/SIL & Faible & Vérifier les types de données, solver settings \\
Accès matériel PIL limité & Moyen & Utiliser PIL simulé (QEMU) ou MCU de développement \\
\bottomrule
\end{tabular}
\end{center}

\end{document}
